\xchapter{总结与展望}{Summary and Outlook}
创新点一（镇守第三章）： 建立了考虑组分重分布与孔隙迁移的快堆MOX燃料微观重构耦合模型。

针对快堆MOX燃料在高功率梯度下的微观演化特征，构建了包含Pu/U组分重分布、氧(O/M)热扩散及孔隙迁移的统一本构模型，修正了传统程序忽略微观结构演化对热物性动态反馈的缺陷，显著提升了高燃耗下物理建模的保真度。

创新点二（镇守第四章）： 提出了一种处理强非线性力学-裂变气体-热耦合效应的隐式数值求解策略。

针对多物理场强耦合计算易发散的难题，开发了基于径向返回算法（Radial Return Mapping）的力学积分求解器，并设计了力学变形-间隙传热-气体释放的全局耦合迭代算法，实现了对燃料元件复杂服役行为的稳定、精确求解。

创新点三（镇守第六章）： 构建了面向数字孪生的基于高保真机理数据驱动的快速预测方法。

为满足数字孪生实时性需求，提出了一种利用高保真模拟数据训练深度神经网络/代理模型的方法。在保持物理模型（创新点1）精度的前提下，实现了对燃料元件关键性能参数的毫秒级实时预报，打通了从离线分析到在线监测的技术瓶颈。

